\documentclass[11pt]{article}
\usepackage[letterpaper,margin=1in]{geometry}
\usepackage{amsmath,amssymb,amsthm,mathtools}
\usepackage{graphicx}
\usepackage{xcolor}
\usepackage[hidelinks]{hyperref}
\usepackage{microtype}

\newtheorem{theorem}{Theorem}
\newtheorem{lemma}[theorem]{Lemma}
\newtheorem{corollary}[theorem]{Corollary}
\newcommand{\Q}{\{-1,1\}}
\newcommand{\E}{\mathbb E}
\newcommand{\N}{\mathbb N}
\newcommand{\R}{\mathbb R}
\newcommand{\ip}[2]{\langle #1,#2\rangle}
\DeclareMathOperator{\supp}{supp}

\title{A finite-dimensional signed $L^1$ uncertainty inequality\\
on the Boolean cube}
\author{Run \texttt{fa\_banach\_001}, lane 18
  (\texttt{agent\_lane\_18})}
\date{12 August 2026}

\begin{document}
\maketitle

\begin{abstract}
Polyanskiy and Samorodnitsky conjectured that their sharp asymptotic
$L^1$ uncertainty estimate for nonnegative low-degree functions on the
Boolean cube remains true for arbitrary real functions.  We prove the
stronger finite-dimensional inequality
\[
 \|f\|_\infty\le
 \sqrt{\frac{\binom{2n}{d}}{\binom{2n}{n}}}
 \sum_{x\in\Q^n}|f(x)|,
 \qquad \deg f\le d.
\]
Its logarithm is
$-n[\log 2-h(d/(2n))]+O(\log n)$, exactly the conjectured exponent.
The proof translates a maximizer to the all-ones vertex, symmetrizes over
coordinates, and uses the exact Christoffel function
$K_d(0,0)=\binom{2n}{d}$ for the squared binomial (Hahn) measure.
\end{abstract}

\section{Source conjecture}

Let $h(t)=-t\log t-(1-t)\log(1-t)$ be binary entropy in natural units.
After proving an $L^2$ uncertainty principle, Polyanskiy and
Samorodnitsky compare it with a result for nonnegative functions and state
the following conjecture on page 20 of the published manuscript
\cite{PS}:

\begin{center}
\includegraphics[width=.91\textwidth]{figures/open_problem_crop.png}
\end{center}

In exact mathematical form, their equation (70) says that, for
$f\ge0$ of Fourier--Walsh degree at most $\delta n$,
\begin{equation}\label{source-estimate}
 \max_{x\in\Q^n}|f(x)|
 \le \exp\!\left(-n\left[\log2-h\left(\frac\delta2\right)\right]
        +o(n)\right)
 \sum_{x\in\Q^n}|f(x)|.
\end{equation}
The immediately following sentence is:
\begin{quote}
``We conjecture that (70) holds for all
$f:\{\pm1\}^n\to\mathbb R$ of degree $\le\delta n$.''
\end{quote}
The result below proves this statement and replaces the unspecified $o(n)$
by an explicit $O(\log n)$ remainder.

\section{Definitions and proof intuition}

Write $\chi_S(x)=\prod_{i\in S}x_i$.  The Fourier--Walsh degree of
$f:\Q^n\to\R$ is the largest $|S|$ for which the coefficient of
$\chi_S$ is nonzero.  All $L^1$ norms in the theorem use counting measure:
\[
 \|f\|_{\ell^1}=\sum_{x\in\Q^n}|f(x)|,
 \qquad \|f\|_\infty=\max_{x\in\Q^n}|f(x)|.
\]

The proof has two independent ideas.  First, translate a point attaining
$\|f\|_\infty$ to $(1,\ldots,1)$ and average over all coordinate
permutations.  Translation preserves Fourier degree, the average keeps the
chosen value, and convexity makes its $\ell^1$ norm no larger.  The averaged
function depends only on the number $k$ of negative coordinates, hence is a
univariate polynomial $P(k)$ of degree at most $d$.

Second, the $k$-distribution of a uniform cube point is binomial.  Squaring
this binomial weight produces a Hahn orthogonality measure.  Its degree-$d$
Christoffel function at the endpoint is exactly $\binom{2n}{d}$; one
Cauchy--Schwarz inequality then bounds $P(0)$ by the binomially weighted
$L^1$ norm of $P$.

\section{The exact theorem}

\begin{theorem}[signed $L^1$ uncertainty]\label{main}
Let $0\le d\le n$.  Every real function $f:\Q^n\to\R$ of
Fourier--Walsh degree at most $d$ satisfies
\begin{equation}\label{finite}
 \boxed{\quad
 \|f\|_\infty\le
 \sqrt{\frac{\binom{2n}{d}}{\binom{2n}{n}}}\,
 \|f\|_{\ell^1}.
 \quad}
\end{equation}
\end{theorem}

\begin{corollary}[the source conjecture]\label{asymptotic}
Uniformly for $0\le d\le n$,
\begin{equation}\label{stirling}
 \sqrt{\frac{\binom{2n}{d}}{\binom{2n}{n}}}
 =\exp\!\left(-n\left[\log2-h\left(\frac{d}{2n}\right)\right]
       +O(\log n)\right).
\end{equation}
Consequently, if $d\le\delta n$, then \eqref{source-estimate} holds for
arbitrary real $f$.
\end{corollary}

\begin{proof}
The uniform estimate
$\log\binom Nk=Nh(k/N)+O(\log N)$ gives \eqref{stirling}.
For $d\le\delta n\le n$, monotonicity of $\binom{2n}{d}$ in this range
completes the deduction from Theorem~\ref{main}.
\end{proof}

\section{The squared-binomial Christoffel identity}

For $0\le k\le n$, put
\[
 \mu_k=2^{-n}\binom nk,
 \qquad
 Z=\sum_{k=0}^n\mu_k^2
   =4^{-n}\binom{2n}{n},
 \qquad
 \nu_k=\frac{\mu_k^2}{Z}
   =\frac{\binom nk^2}{\binom{2n}{n}}.
\]
The identity for $Z$ is Vandermonde's convolution.

For a probability measure $\nu$ on $\{0,\ldots,n\}$, let
\[
 K_d^\nu(0,0)=
 \sup_{0\ne P,\ \deg P\le d}
 \frac{|P(0)|^2}{\sum_{k=0}^n\nu_k|P(k)|^2}
\]
be its degree-$d$ Christoffel function at zero.

\begin{lemma}[endpoint Hahn identity]\label{hahn}
For the squared-binomial probability measure above,
\[
 K_d^\nu(0,0)=\binom{2n}{d},\qquad 0\le d\le n.
\]
\end{lemma}

\begin{proof}
We record the short normalization calculation so that no asymptotic
orthogonal-polynomial result is hidden in the proof.  With $(a)_r$ denoting
the rising factorial, define the terminating Hahn polynomials
\[
 H_j(k)={}_3F_2\!\left(
 \begin{matrix}-j,\ j-2n-1,\ -k\\-n,\ -n\end{matrix};1\right)
 =\sum_{r=0}^j
 \frac{(-j)_r(j-2n-1)_r(-k)_r}{(-n)_r^2r!}.
\]
They have degree $j$ and $H_j(0)=1$.  Expanding the terminating sums and
applying Vandermonde's convolution (equivalently, the standard finite Hahn
orthogonality calculation \cite[Sec.~9.5]{KLS}) gives
\begin{equation}\label{orthogonality}
 \sum_{k=0}^n\nu_k H_i(k)H_j(k)
 =\frac{\mathbf1_{\{i=j\}}}
 {\binom{2n}{j}-\binom{2n}{j-1}},
\end{equation}
where $\binom{2n}{-1}=0$.

Thus $H_0,\ldots,H_d$ are an orthogonal basis of the polynomials of degree
at most $d$.  The reproducing-kernel formula and $H_j(0)=1$ yield
\[
 K_d^\nu(0,0)
 =\sum_{j=0}^d\frac{|H_j(0)|^2}{\|H_j\|_{L^2(\nu)}^2}
 =\sum_{j=0}^d
 \left[\binom{2n}{j}-\binom{2n}{j-1}\right]
 =\binom{2n}{d}.
\]
This also proves the claimed normalization directly.  The accompanying
exact-arithmetic script expands \eqref{orthogonality} for all degrees through
$n=12$ as an independent check.
\end{proof}

\begin{lemma}[radial $L^1$ evaluation]\label{radial}
Every real polynomial $P$ of degree at most $d\le n$ satisfies
\begin{equation}\label{radial-bound}
 |P(0)|\le
 \sqrt{\frac{\binom{2n}{d}}{Z}}
 \sum_{k=0}^n\mu_k|P(k)|.
\end{equation}
\end{lemma}

\begin{proof}
Set $A=\sum_k\mu_k|P(k)|$.  Since all summands are nonnegative,
\[
 A^2\ge\sum_{k=0}^n\mu_k^2|P(k)|^2
 =Z\|P\|_{L^2(\nu)}^2.
\]
The definition of the Christoffel function and Lemma~\ref{hahn} give
\[
 |P(0)|^2\le\binom{2n}{d}\|P\|_{L^2(\nu)}^2
 \le\frac{\binom{2n}{d}}{Z}A^2.
\]
Taking square roots proves \eqref{radial-bound}.
\end{proof}

\section{Proof of the theorem}

\begin{proof}[Proof of Theorem~\ref{main}]
Choose $x^0\in\Q^n$ with $|f(x^0)|=\|f\|_\infty$ and multiply $f$ by a
sign.  The cube translation
\[
 g(x)=f(x^0_1x_1,\ldots,x^0_nx_n)
\]
has the same Fourier degree and the same $\ell^1$ norm, and
$g(1,\ldots,1)=\|f\|_\infty$.

Average $g$ over the symmetric group:
\[
 \bar g(x)=\frac1{n!}\sum_{\pi\in S_n}
 g(x_{\pi(1)},\ldots,x_{\pi(n)}).
\]
Then $\deg\bar g\le d$, the all-ones vertex is fixed, and Jensen's
inequality followed by permutation invariance of counting measure gives
\begin{equation}\label{contraction}
 \bar g(1,\ldots,1)=\|f\|_\infty,
 \qquad
 \|\bar g\|_{\ell^1}\le\|g\|_{\ell^1}=\|f\|_{\ell^1}.
\end{equation}

A symmetric Fourier--Walsh polynomial of degree at most $d$ is a linear
combination of the first $d$ Krawtchouk polynomials.  Hence
$\bar g(x)=P(k)$, where $k=|\{i:x_i=-1\}|$ and $P$ is a univariate
polynomial of degree at most $d$.  In particular,
\[
 \sum_{k=0}^n\mu_k|P(k)|
 =2^{-n}\|\bar g\|_{\ell^1},
 \qquad P(0)=\|f\|_\infty.
\]
Lemma~\ref{radial}, the value
$Z=4^{-n}\binom{2n}{n}$, and \eqref{contraction} now give
\begin{align*}
 \|f\|_\infty
 &\le
 \sqrt{\frac{\binom{2n}{d}}
 {4^{-n}\binom{2n}{n}}}\;2^{-n}
 \|\bar g\|_{\ell^1}\\
 &\le
 \sqrt{\frac{\binom{2n}{d}}{\binom{2n}{n}}}
 \|f\|_{\ell^1},
\end{align*}
which is \eqref{finite}.
\end{proof}

\section{Verification, novelty, and limitations}

There are several quick checks on the exact constant.  If $d=n$, it is one
and \eqref{finite} is the trivial inequality $\|f\|_\infty\le\|f\|_1$.
If $d=0$, it differs from the sharp constant $2^{-n}$ only by the polynomial
factor $4^{n/2}/\sqrt{\binom{2n}{n}}=O(n^{1/4})$, so it has the correct
exponential scale.  Exact symbolic computation verifies the Hahn
orthogonality, reciprocal norms, and telescoping identity for every
$1\le n\le12$; see \texttt{verify\_hahn\_identity.py}.

Targeted searches of the source, its published version, the run registry,
Boolean-cube Nikolskii inequalities, the exact binomial-ratio constant, and
squared-binomial Christoffel functions found no prior statement of
Theorem~\ref{main}.  The source TeX contains an appendix-only working note
mentioning a signed permutation-invariant estimate, and Samorodnitsky's older
preprint \cite{Sam} develops the relevant squared-binomial/Hahn machinery.
Those are consistent with Lemma~\ref{radial}; neither located source states the
translation--symmetrization reduction or the finite inequality
\eqref{finite}.  Literature novelty is nevertheless the main point requiring
expert review.

The theorem resolves only the signed $L^1$ conjecture following equation
(70).  It does not solve the source's separate conjecture that Hamming balls
are asymptotically worst for support-constrained hypercontractivity, nor its
more structured uncertainty questions.

\paragraph{Human review recommendation.}
Before dissemination, verify the normalization in \eqref{orthogonality}, the
degree preservation in the radial reduction, and novelty against the
literature on Hahn Christoffel functions and discrete Nikolskii inequalities.

\begin{thebibliography}{9}
\bibitem{PS}
Y.~Polyanskiy and A.~Samorodnitsky,
\emph{Improved log-Sobolev inequalities, hypercontractivity and uncertainty
principle on the hypercube}, J. Funct. Anal. 277 (2019), 108280;
arXiv:1606.07491.

\bibitem{KLS}
R.~Koekoek, P.~A. Lesky, and R.~F. Swarttouw,
\emph{Hypergeometric Orthogonal Polynomials and Their $q$-Analogues},
Springer, 2010.

\bibitem{Sam}
A.~Samorodnitsky,
\emph{Extremal properties of solutions for Delsarte's linear program},
preliminary manuscript, 1998.
\end{thebibliography}

\end{document}
